\section{Considerações Finais}

Este trabalho apresentou uma abordagem baseada em Mineração de Dados para a caracterização de perfis acadêmicos de estudantes de graduação, tendo como objetivo contribuir para o desenvolvimento de uma ferramenta capaz de oferecer suporte personalizado à trajetória acadêmica. A proposta considera informações relacionadas ao desempenho dos estudantes e ao histórico de aprovações e reprovações como elementos para a identificação de padrões acadêmicos.

No estágio desenvolvido nesta pesquisa, foi implementada uma arquitetura composta por um chatbot integrado ao Telegram, responsável pela interação com o estudante e pelo recebimento dos históricos acadêmicos, além de uma etapa de extra\c{c}\~ao e tratamento dos dados e da
aplica\c{c}\~ao dos algoritmos K-Means, para agrupamento dos estudantes
segundo caracter\'isticas de desempenho, e KNN, para classifica\c{c}\~ao
por eixo de forma\c{c}\~ao e por faixa de desempenho acad\^emico geral.

A pesquisa também avaliou a possibilidade de incorporar Grandes Modelos de Linguagem à arquitetura do tutor, considerando seu potencial para proporcionar interações mais naturais e personalizadas. Entretanto, as restrições financeiras e operacionais identificadas durante a análise de viabilidade inviabilizaram sua utilização na configuração atual do projeto. Dessa forma, optou-se pela continuidade da abordagem baseada em técnicas de Mineração de Dados, mantendo a integração com LLMs como possibilidade para etapas futuras.

Os experimentos realizados mostram a ideia proposta e como executá-la. Os resultados obtidos constituem uma etapa inicial para o desenvolvimento de um sistema de recomendação acadêmica. A identificação de perfis semelhantes pode futuramente ser utilizada em conjunto com informações sobre pré-requisitos, disciplinas pendentes e trajetórias de estudantes concluintes para auxiliar na definição de caminhos acadêmicos personalizados.

Assim, o trabalho estabelece uma base para a evolução de um tutor inteligente orientado por dados, cuja efetividade poderá ser avaliada em etapas posteriores a partir de dados reais e de interações com estudantes da UFRRJ.

\subsection{Trabalhos Futuros}
\label{sec:trabalhos_futuros}

Como continuidade deste trabalho, pretende-se, primeiramente, verificar o funcionamento da abordagem desenvolvida e testá-la com dados reais, a partir dos históricos escolares a serem disponibilizados pela COTIC. Isso viabilizará a consolidação de um conjunto de dados representativo da comunidade discente de Ciência da Computação da UFRRJ e a replicação dos experimentos. Vale destacar que, à medida que os estudantes utilizarem o sistema, mais dados serão inseridos na base, e a tendência é que os algoritmos melhorem sua taxa de acerto.

Outra frente é a análise dos históricos acadêmicos já obtidos pelo grupo. Um código em Python já desenvolvido extrai de cada histórico a porcentagem de integralização do curso; adaptando-o para processar todos os históricos em lote, será possível construir uma curva que indique em quais pontos do curso os estudantes mais ficam retidos. Essa informação poderá servir como entrada para os algoritmos de agrupamento e de recomendação.

Na sequência, prevê-se a finalização do módulo de recomendação acadêmica, integrando os perfis mapeados por agrupamento (clustering) às matrizes de pendências curriculares, encadeamento de pré-requisitos e trajetórias históricas de concluintes. Essa integração permitirá investigar algoritmos para a geração de planos de estudo personalizados, avaliando sua eficácia no fluxo de integralização curricular.

Adicionalmente, projeta-se a incorporação de Grandes Modelos de Linguagem como camada conversacional do tutor inteligente. Nessa arquitetura híbrida, os modelos de Mineração de Dados Educacionais atuarão no processamento analítico dos dados, enquanto a LLM exercerá o papel de sintetizar e contextualizar os diagnósticos para o usuário. A exequibilidade dessa etapa permanece condicionada à captação de recursos orçamentários ou à alocação de infraestrutura computacional capaz de suportar acessos simultâneos.

Por fim, pretende-se conduzir uma avaliação empírica com estudantes da instituição, aferindo métricas de usabilidade, aceitação da tecnologia e assertividade das recomendações geradas. Em longo prazo, mediante a consolidação de dados longitudinais, investigar-se-á o impacto do sistema sobre indicadores acadêmicos chave, como taxas de retenção, tempo de integralização e índices de retenção e evasão.